\section{Project Roadmap}\label{sec:sim-roadmap}

\subsection*{Scope of this chapter}
\addcontentsline{toc}{subsection}{Scope of this chapter}

This chapter is the roadmap for the \textbf{Simula training-data
framework} project. It is a project-level artefact, not a programme
plan. Other SAP OSS Finance projects (Trial Balance, Arabic AP,
Regulations) run on the same platform, workspaces, and backends, but
each owns its own roadmap, business case, and go/no-go decision.
Simula is a \emph{supplier} to those projects: it generates reference
training corpora they consume, so its critical path is sequenced
before theirs, but its sign-off is independent.

\paragraph{Sequence, not dates.}
The milestones and dependency graph below express \emph{order} and
\emph{prerequisite} relationships, not calendar dates. This is
deliberate. Implementation teams must know what depends on what, and
where the schedule can absorb slippage (the Should/Could items); a
date column would harden targets into commitments they were never
meant to be. Dates, when set, live in the project steering packet and
the ticketing system, not in this specification.

\subsection{MoSCoW classification}\label{sec:sim-moscow}

Every milestone below carries one of four priorities. The mapping is
contractual: a milestone's priority determines what happens when it is
dropped, de-scoped, or slips.

\begin{table}[htbp]
\centering
\caption{MoSCoW prioritisation (shared across the four SAP OSS Finance
projects).}
\label{tab:sim-moscow}
\footnotesize
\begin{tabularx}{\linewidth}{@{}l l X@{}}
\toprule
\textbf{Priority} & \textbf{Meaning} & \textbf{Allowance} \\ \midrule
\textbf{Must} (\textbf{M}) & Non-negotiable for v1.2. Without this milestone, the release does not ship. & None. Missing a Must halts the release; see the kill-switch in \cref{sec:sim-kill-switch}. \\
\textbf{Should} (\textbf{S}) & Important but not vital. Ship without it only with sponsor sign-off; document the deviation in the business-case refresh. & May slip by one consumer project; two consecutive slips escalate to GFAIC. \\
\textbf{Could} (\textbf{C}) & Desirable enhancement. Included if Must and Should items are on track. & May be dropped to v1.3 at project steering's discretion. \\
\textbf{Won't} (\textbf{W}) & Explicitly out of v1.2 scope. Listed so the boundary is legible. & Not eligible for v1.2 planning at all. \\
\bottomrule
\end{tabularx}
\end{table}

\noindent
The critical path of this project is the chain of Must items in
\cref{tab:sim-milestones}; the Should and Could items are where the
schedule and scope can absorb delivery risk without a release slip.

\subsection{Dependency graph}\label{sec:sim-dependency-graph}

\begin{figure}[htbp]
\centering
\begin{tikzpicture}[
  node distance=5mm and 10mm,
  must/.style={draw,rounded corners=2pt,align=center,font=\scriptsize,
               fill=sapgold!20,draw=sapgold!80,
               minimum height=7mm,minimum width=32mm},
  should/.style={draw,rounded corners=2pt,align=center,font=\scriptsize,
               fill=sapblue!15,draw=sapblue!60,
               minimum height=7mm,minimum width=32mm},
  could/.style={draw,rounded corners=2pt,align=center,font=\scriptsize,
               fill=hanagreen!15,draw=hanagreen!60,dotted,
               minimum height=7mm,minimum width=32mm},
  dep/.style={-{Latex[length=2mm]},thick},
  soft/.style={-{Latex[length=2mm]},thick,dashed}
]
  \node[must] (m01) {M01 SIM-ENG-001 \\ \scriptsize extraction pipeline};
  \node[must,below=of m01] (m02) {M02 SIM-ENG-002 \\ \scriptsize taxonomy engine};
  \node[must,below=of m02] (m03) {M03 SIM-ENG-003 \\ \scriptsize generation engine};
  \node[must,right=14mm of m02] (m04) {M04 SIM-ENV-001 \\ \scriptsize environment CLI};
  \node[must,right=14mm of m04] (m05) {M05 SIM-HANA-001 \\ \scriptsize spatial policy};
  \node[must,below=of m05] (m06) {M06 SIM-RUN-TB \\ \scriptsize TB reference run};
  \node[must,below=of m06] (m07) {M07 SIM-RELEASE-001 \\ \scriptsize v1.2 tag};
  \node[should,right=14mm of m07] (s01) {S01 SIM-RUN-AP \\ \scriptsize AP corpus};
  \node[should,above=of s01] (s02) {S02 SIM-RUN-REG \\ \scriptsize Reg corpus};
  \node[could,above=of s02] (c01) {C01 SIM-V13-EXPLORE \\ \scriptsize v1.3 scope};

  \draw[dep] (m01) -- (m02);
  \draw[dep] (m02) -- (m03);
  \draw[dep] (m03) -- (m04);
  \draw[dep] (m04) -- (m05);
  \draw[dep] (m05) -- (m06);
  \draw[dep] (m06) -- (m07);
  \draw[dep] (m07) -- (s01);
  \draw[dep] (s01) -- (s02);
  \draw[soft] (s02) -- (c01);
\end{tikzpicture}
\caption{Simula dependency graph. Amber nodes are Must (critical
path); blue are Should (consumer-corpus allowance); dotted green are
Could (discretionary). Solid arrows are hard prerequisites; dashed
arrows indicate conditional starts (the target is unlocked when the
upstream completes \emph{and} the v1.3 trigger in \cref{sec:sim-v13-trigger}
fires).}
\label{fig:sim-dependency-graph}
\end{figure}

\subsection{Project milestones}\label{sec:sim-milestones}

\begin{table}[htbp]
\centering
\caption{Simula v1.2 project milestones in dependency order.
Priority: M=Must (critical path), S=Should, C=Could, W=Won't.}
\label{tab:sim-milestones}
\footnotesize
\begin{tabularx}{\linewidth}{@{}l l l l X@{}}
\toprule
\textbf{ID} & \textbf{Prereq} & \textbf{Pri.} & \textbf{Ticket} & \textbf{Milestone} \\ \midrule
M01 & ---      & M & SIM-ENG-001       & Extraction pipeline passes HANA-direct fixtures in the testing chapter. \\
M02 & M01      & M & SIM-ENG-002       & Taxonomy engine classifies reference corpus with accuracy above Beta threshold. \\
M03 & M02      & M & SIM-ENG-003       & Generation engine produces schema-valid examples against the unified registry. \\
M04 & M03      & M & SIM-ENV-001       & Environment CLI enables one-command reproduce of any reference run. \\
M05 & M04      & M & SIM-HANA-001      & HANA spatial policy (\texttt{spatial\_unknown\_srid\_policy}) enforced; strict mode passes unit tests. \\
M06 & M05      & M & SIM-RUN-TB        & Trial Balance reference run reproducible end-to-end (commit \texttt{b8d40a1}, tag \texttt{simula-v1.2-worked-example}, seed \texttt{20251115}). \\
M07 & M06      & M & SIM-RELEASE-001   & v1.2 release tagged, archived, and the reproducibility manifest frozen. \\
M08 & M07      & M & SIM-BC-REFRESH    & Business-case refresh (first cycle after release; then annual). \\
S01 & M07      & S & SIM-RUN-AP        & Arabic AP training corpus generated from the same pipeline. \\
S02 & S01      & S & SIM-RUN-REG       & Regulations monitoring corpus generated from the same pipeline. \\
S03 & M07      & S & SIM-CONSUMER-DOC  & Consumer integration runbook published (how Trial Balance, AP, and Regulations call Simula). \\
C01 & S02 + v1.3 trigger (\S\ref{sec:sim-v13-trigger}) & C & SIM-V13-EXPLORE & v1.3 scope kick-off (multi-tenant isolation, additional taxonomies). \\
C02 & M07      & C & SIM-TUNE          & Quarterly prompt and classifier retraining from accept/edit telemetry. \\
W01 & ---      & W & SIM-V12-EXTERNAL  & External (non-SAP-OSS) customer access is \emph{not} in v1.2 scope. \\
W02 & ---      & W & SIM-V12-REALTIME  & Real-time incremental generation is \emph{not} in v1.2 scope; batch only. \\
\bottomrule
\end{tabularx}
\end{table}

\noindent
The Must chain (M01 $\to$ M02 $\to$ M03 $\to$ M04 $\to$ M05 $\to$ M06
$\to$ M07 $\to$ M08) is the critical path. The Should items
(S01--S03) are the primary \emph{schedule allowance}: consumer corpus
generations can slip without affecting Simula v1.2 sign-off. The Could
items are the \emph{scope allowance}: they may be cut to v1.3 without
renegotiation.

\subsection{Launch tiers and exit gates}\label{sec:sim-launch-tiers}

The project walks the four-tier ladder shared across SAP OSS Finance
projects. A tier is exited only when the named exit gate is satisfied;
slippage requires sponsor sign-off.

\begin{table}[htbp]
\centering
\caption{Simula launch tiers. The exit gate defines the MoSCoW items
that must be closed to exit the tier.}
\label{tab:sim-launch-tiers}
\footnotesize
\begin{tabularx}{\linewidth}{@{}l l X@{}}
\toprule
\textbf{Tier} & \textbf{Scope} & \textbf{Exit gate (references MoSCoW IDs)} \\ \midrule
Alpha & Engineering only, DEV + UAT, fixture corpora only. & M01, M02, M03 closed; schema-registry validation passes on the generated examples; reproducibility confirmed on two consecutive runs. \\
Beta  & One consumer project (Trial Balance) drawing from Simula-generated corpora. & M04, M05, M06 closed; environment CLI reproduces the TB run end-to-end; HANA spatial strict mode green; no open SEV1/SEV2. \\
GA (limited) & Two consumer projects (TB plus Arabic AP). & M07, M08 closed; S01 closed; consumer-integration runbook (S03) exercised end-to-end; kill-switch green. \\
Full rollout & All three current consumers (TB, AP, Regulations). & S02 closed; two quarterly reviews sustained; business-case refresh complete; multi-consumer regression suite green. \\
\bottomrule
\end{tabularx}
\end{table}

\subsection{Kill-switch conditions}\label{sec:sim-kill-switch}

The kill-switch is a pre-committed rulebook: a tripped condition enters
the action column without further debate. GFAIC and sponsor are
notified within 24 hours. The conditions are evaluated continuously
(on every cycle) and do not carry dates.

\begin{table}[htbp]
\centering
\caption{Simula kill-switch.}
\label{tab:sim-kill-switch}
\footnotesize
\begin{tabularx}{\linewidth}{@{}l l X@{}}
\toprule
\textbf{Condition} & \textbf{Evaluation} & \textbf{Action} \\ \midrule
Reference run non-reproducible across two consecutive seeds (output diff $>$\,0) & per release & \textbf{Halt release}; investigate non-determinism; do not ship v1.2 until reproducibility is restored. \\
Schema-registry validation fails on $>$\,0.1\,\% of generated examples & per run & \textbf{Halt consumer delivery}; fix generation engine; re-run. \\
HANA spatial strict mode detects unknown SRID in production path & continuous & \textbf{Quarantine run}; incident-review before releasing corpus to consumer project. \\
Any consumer project (TB, AP, Regulations) declares red on a Simula-supplied corpus & continuous & \textbf{Pause Should items} S01--S03; remediate before next consumer wave. \\
Shared-platform dependency in \cref{tab:sim-shared-deps} declared red by its owner project & continuous & \textbf{Pause downstream Must}; re-plan at next steering. \\
\bottomrule
\end{tabularx}
\end{table}

\subsection{Shared-platform dependencies}\label{sec:sim-shared-deps}

This project depends on three shared-platform components. Each is
tracked here only as a dependency, not a commitment this project can
unilaterally adjust. Sequence is expressed in terms of the Must item
that first consumes the dependency, not in calendar time.

\begin{table}[htbp]
\centering
\caption{Shared-platform dependencies.}
\label{tab:sim-shared-deps}
\footnotesize
\begin{tabularx}{\linewidth}{@{}l l l X@{}}
\toprule
\textbf{Component} & \textbf{Owner} & \textbf{First consumed by} & \textbf{Impact if late} \\ \midrule
Unified schema registry (\path{docs/schema/}) & Shared infra & M03 & Already \texttt{migration.status=complete}; no remaining risk. \\
vLLM TurboQuant hosting (\texttt{vllm-only} policy) & Shared platform & M03 & Graceful-degradation ladder (\cref{sec:degradation-ladder} of the Regulations spec) provides fallback; batch-generation slip tolerable. \\
HANA Cloud spatial extensions & Shared infra & M05 & No substitute; slip blocks spatial-policy certification and downstream consumer runs. \\
\bottomrule
\end{tabularx}
\end{table}

\noindent
Shared-platform dependencies are surfaced at monthly steering. If any
upstream owner declares \textbf{amber} or \textbf{red} on a component
in \cref{tab:sim-shared-deps}, the affected downstream Must milestone
is re-planned at the following steering.

\subsection{v1.3 expansion trigger}\label{sec:sim-v13-trigger}

The Could item C01 (SIM-V13-EXPLORE) is unlocked only when all of the
following are observed after two consecutive consumer corpus
generations under v1.2:

\begin{itemize}
  \item aggregate consumer request volume for corpora outside the
        current taxonomy exceeds 20\,\% of total Simula runs;
  \item reproducibility and schema-validation rates sustained at or
        above release thresholds;
  \item business-case refresh (M08) explicitly budgets the v1.3
        programme.
\end{itemize}

\noindent
If any of the three conditions fails, C01 remains in the Could column
and is re-evaluated at the next business-case refresh.

\subsection{Governance and business-case cadence}\label{sec:sim-governance}

\paragraph{Project governance.}
The Simula project reports monthly to a project steering (Studio AI
Lead, Training Data PM, Pipeline Engineer) and quarterly to the
cross-project GFAIC. Monthly steering accepts or rejects amber/red
signals on Must items; GFAIC reviews compliance posture and any
cross-project dependencies (especially consumer-project risk).

\paragraph{Business-case refresh (M08).}
The business case that authorises this project is refreshed on the
following rhythm:

\begin{itemize}
  \item \textbf{Mandatory refresh}: at the first cycle after release,
        and annually thereafter (ticket SIM-BC-REFRESH, priority M in
        \cref{tab:sim-milestones}). Refresh packet includes consumer
        demand assumptions, unit economics (\cref{sec:sim-unit-economics}),
        GPU cost realised vs.\ forecast, and a deviation statement
        against the baseline.
  \item \textbf{Trigger-based refresh}: any of
        \begin{itemize}
          \item consumer request volume changes by $\ge\,20$\,\% from
                forecast (either direction);
          \item per-run generation cost changes by $\ge\,30$\,\% (GPU
                price, vLLM licensing, HANA I/O);
          \item a consumer project's regulatory obligation expands
                Simula's in-scope taxonomy;
          \item any kill-switch in \cref{tab:sim-kill-switch} is
                tripped.
        \end{itemize}
  \item \textbf{Owner}: Studio AI Lead. Approval routed through the
        project steering and filed at
        \path{docs/simula/business-case/refresh-history.yaml}.
\end{itemize}
